\documentclass[11pt]{article}
\usepackage{amsmath,amssymb,amsthm,geometry}\geometry{margin=1in}
\newtheorem{theorem}{Theorem}\newtheorem{proposition}{Proposition}
\title{Coherent primitive-occurrence instance libraries}
\date{17 November 2026}\begin{document}\maketitle
\paragraph{Construction and scope.}
This constructs the fresh-instance supply required by the support-separated atomic polynomial in \texttt{\detokenize{Entry_Sponsored_Atomic_Progressive_Polynomials.tex}}. It is not an inference that every fixed-name graph already contains arbitrarily many disjoint copies. Instances are indexed deterministically by retained primitive request occurrences and event ancestry. No fresh-name probability distribution or new paid event is introduced. The indexing is extra retained interface data; its eventual forgetting is a specified map, not generally an isomorphism of opaque states.

\section{Local reusable data and allocation indices}
A reusable component type $\tau$ has a small local native library $G_\tau$: states, all command schemas, finite raw initializer/paid rows or declared families, actual coefficients, full terminal and writer grades, selected-entry and poststate functions, old-port inclusions and the required world actions. Its primitive actor and material-request occurrences have named local slots, and every native field uses those slots consistently. All old local occurrences persist along the operations in the admitted branch. A raw-dependent birth profile may vary; its finite slot universe for a fixed row is the finite union of all raw request addresses, with absent sentinels, not merely the positive support. The all-raw finite-initialization guard supplies this rowwise universe as in the owner-aware activation theorem.

Names are typed symbols, not observed values. Let $L_\tau^{\rm act}$ and $L_\tau^{\rm mat}$ be disjoint sets of local actor and material slots across the library. They may be infinite across all inputs, although each realized state/row uses finitely many. A material slot exists even when its initializer selects a state with no actor or output. These local occurrence presentations are part of the source interface; an arbitrary opaque system without any such occurrence presentation is not silently assigned one.

A run has a retained unique run token $R$, and each root occurrence identifier $r$ is qualified by that token. Distinct runs combined in one library use distinct tokens or a supplied injective run reindexing. Retain an event index $n=(R,k)$ with ordinal $k\in\mathbb N$ for each actual atomic input occurrence, including paid identity-world events and declared pure input events. The ordinal is global across the relevant combined run, not a separate counter independently reused by each root. Run qualification applies even when a pure-input sponsor is not itself rooted at $r$. The ordinal merely enumerates existing atoms: it neither creates a paid tick nor subdivides a scope. Silent initializer descendants within an atom use request ancestry, not extra event ordinals. Initial installation uses the distinct tag $\mathrm{init}$.

A new invocation request is identified by its primitive producing sponsor $o$, its event occurrence $n$ (or initial root tag), its formal initializer/port ancestry path $p$, and its component type $\tau$. The path $p$ uses primitive request/port slots and silent ancestor slots, never content grouping wrappers. The complete request signature must distinguish every simultaneous invocation in this way. Request types and presence are determined by the declared predecessor values; an absent slot creates no instance. Repeated observed labels do not identify request slots.

Use a free typed term algebra of namespaces. Root instances have namespaces $\operatorname{Root}(r,p,\tau)$; paid or pure-input births have
\[
          a=\operatorname{Birth}(n,o,p,\tau).              \tag{1}
\]
Actor and material names in this instance are respectively
\[
       \operatorname{Actor}(a,l),\qquad\operatorname{Material}(a,m).
                                                               \tag{2}
\]
The two constructors are disjoint. Primitive request records that belong to the enclosing activation rather than the internal component use $\operatorname{Request}(n,o,p)$, a third disjoint sort. If a formal local slot is itself hierarchical, normalize it by the supplied primitive ancestry representation; nested content parentheses are not added to (1). Existing actor names remain unchanged when an event is processed.

\section{The library of actual copies}
For each namespace $a$ of type $\tau$, construct $G_\tau^a$ by a literal tagged copy of every state and raw witness and by substituting (2) for local occurrence names throughout every interface field. Copy every initializer/command law unchanged on its bijective raw carrier. Copy actual context and content coefficient objects and maps, not their scalar values. The entire unused terminal grade remains the same full $C_\tau$ with its actual maps; formal addressed writer schemas are renamed at every schema, including unused ones. Initial selected states, paid successors, boundary maps, passive responses and old-port inclusions are defined by conjugating the original functions with the copy bijections.

The reusable instance library is the disjoint indexed family $(G_\tau^a)_a$. It is not a stochastic coproduct over $a$. An actual request with index (1) selects that one copy before applying its unchanged initializer raw family. Even if infinitely many copies exist in the library, each finite request diagram activates finitely many and introduces no new random coordinate for choosing a name.

\begin{theorem}[Actual copy and deterministic fresh installation]
At a fixed supplied namespace, the copy bijections give a native isomorphism to the original local occurrence presentation after the declared name reindexing. Deterministic allocation (1)--(2) produces disjoint actor and material supports for all distinct invocations in a finite same-event activation diagram, disjoint also from previously installed instances. It retains every raw zero/absent branch and all actual native fields.
\end{theorem}
\begin{proof}
Conjugating every function gives the initial, successor, world and port squares. Raw bijections push every admitted law exactly; actual coefficient maps are the copied original maps, with identities or canonical reindexing isomorphisms. The full-grade and all-schema fields were copied on their complete domains. Thus no verification is restricted to positive raw support.

Two same-event request indices differ in sponsor, primitive request path or type, hence give distinct namespace terms. Two requests from distinct atomic occurrences differ in $n$. Initial namespaces use a distinct constructor. A descendant's formal ancestry extends its generating path and therefore cannot collide with that ancestor or a sibling. The actor/material/request constructors distinguish different material roles. Thus fresh support follows from injectivity of free term constructors, not from random name search. All-raw induction through the finite request graph allocates even counterfactual zero branches consistently; absence creates only the declared sentinel and no actor/material instance. Old names are never rewritten.\end{proof}

The inverse copy map is a reindexing isomorphism, not a strict fixed-global-name graph map. Once the instances have been supplied on a common named signature, the polynomial graph theorem acts on them using strict name-preserving maps. This distinction prevents the copy construction from contradicting that theorem's narrower map class.

\section{Regrouping, unit and event-history coherence}
\begin{theorem}[Primitive-index coherence]
Suppose two activation recipes retain the same primitive request forest, producing actors, actual event sequence and complete raw assignment, and differ only by content grouping. Their allocated instances, raw selected states and actorless material records agree on the common namespace presentation. The comparisons commute with $\eta$, $\mu$, every finite content associator and their pentagon, and with chronological concatenation of retained event histories.
\end{theorem}
\begin{proof}
Initial root indices agree. At each atomic occurrence, all old sponsor names agree by induction. The next names depend only on that occurrence's $n$, the same primitive sponsor, the same formal request ancestry and type. Hence every new namespace agrees, including entries which already expose downstream requests. Silent descendants extend the same ancestry path and keep their original primitive producing sponsor; no content wrapper contributes a namespace. This proves equality on every raw assignment, including zero and unreachable ones.

Consequently the polynomial multiplication simply rebrackets already allocated tuples. It does not prefix or change global IDs. A unary content unit contributes no actor or material request and leaves all allocation terms unchanged. Four regrouping routes have the same primitive request forest, so their maps agree term by term. Native fields agree through the copy conjugations; actual owner tensors and scope product legs use the same named factors once. Chronological concatenation preserves earlier names and appends the later actual occurrence indices. If suffix histories are stored with relative counters, their names are terms with formal boundary-name parameters. Instantiating after a prefix of length $k$ replaces every locally generated ordinal $m$ by $k+m$ with the same run token, recursively inside sponsor names and any name-bearing ancestry terms as well as at the outer birth node. Boundary parameters are substituted by the already allocated prefix names, which are not shifted again. When two suffixes are composed, substitute the first suffix's name terms for the second's boundary parameters before final instantiation; do not freeze them prematurely as imported global names. Structural recursion on these terms, together with $(k+m)+n=k+(m+n)$ for the locally generated counters, proves the two concatenation routes agree. This is recursive address substitution, not merely offsetting the outermost event label. No event is inserted or erased by this indexing.\end{proof}

If an implementation initially stores nested block addresses, its required comparison is the map from those addresses to their primitive sponsor/request occurrence terms. That map must be proved bijective on the material occurrences of the recipe, including actorless entries, and must carry the actual native fields. It is not sufficient to erase parentheses in an arbitrary opaque address representation. The theorem gives the canonical common target when this primitive occurrence decoder is supplied.

\section{Native maps and the exact index boundary}
A native component map preserving local occurrence slots and types extends copywise at every common namespace: state/raw maps are the copied original maps, and global names in (2) are unchanged. Its initial and paid law pushforwards, whole schema maps and coefficient/body squares remain the original ones under conjugation. Whole-joint owner/request admission is still checked on the same common raw maps; allocation introduces no independent sampling and no hidden-coordinate-dependent witness permutation.

If a component map renames or identifies local occurrences, the induced map has instead the explicit name map
$\operatorname{Actor}(a,l)\mapsto\operatorname{Actor}(a',\rho(l))$ and the corresponding material/request map. Such a comparison is not automatically a strict fixed-name graph map. A noninjective name map can destroy disjointness of simultaneous instances and therefore needs either a restricted map class, a proved disjointness-preserving address comparison, or a different base category with that reindexing built in. Similarly a temporal map erasing or merging atomic occurrences does not preserve (1) without an additional event-index map and its compatibility equations. Paid-length-preserving native history comparisons avoid that particular problem; they do not automatically solve every local occurrence contraction.

\begin{proposition}[No uniform reuse inside an unchanged finite-name presentation]
A fixed-name graph containing only one nonempty supported instance cannot realize two simultaneous disjoint copies merely by reusing that vertex. Nor can a stateless allocation depending only on the reusable type assign distinct names to two simultaneous requests of that type.
\end{proposition}
\begin{proof}
Both copies would contain the same fixed primitive name, contradicting support separation. A function of the same type alone returns the same namespace on both requests. Request-occurrence or freshness-environment information is therefore necessary.\end{proof}

The constructed library supplies that missing information as an external deterministic index. If the original opaque state does not retain an event occurrence counter or equivalent primitive request identity, adding the allocation index is a change of state/input presentation. Forgetting it is a native observation when all required fields factor, but need not have an inverse: two occurrences with the same opaque state can carry different installation histories and hence different fresh names. The claim is rowwise native copying and coherent indexed activation, not an unproved isomorphism between an arbitrary opaque system and its history-enriched realization.

Thus reusable instances are now concretely supplied for the occurrence-presented, event-indexed branch. The receiver action may consume those actual copies and the support-separated polynomial on their common signature. It must still verify its full filled-signature law/material transport; deterministic fresh allocation neither chooses a coupling nor manufactures that receiver action.
\end{document}
